\subsection{Official DGPPO With Task Adaptation and Controlled Budgets}\label{subsec:external}
An independent learning reference uses the public implementation of DGPPO by Zhang et al.\cite{zhang2025dgppo}, retaining its graph-transformer policy, task value network, decomposed constraint value network, discrete-barrier advantage, and proximal policy update. We use the upstream nonrecurrent option and a separate environment adapter replacing the original second-order dynamics with our single integrator. Signed collision, spectral-connectivity, and workspace costs enter its constraint interface. The original training-cost margin transformation is retained; it never modifies physical distances, integrated states, or evaluation audits.

Both methods use the original lattice task generator over $N\in\{6,\ldots,10\}$. DGPPO uses finite banks of $4096$ tasks per scale, with size-homogeneous batches drawn according to accepted sample frequencies; EGP-NPA uses its existing online task stream. DGPPO rollout windows contain $500$ steps, with $480000$ environment transitions per replicate counting both stochastic collection and deterministic collection inside the update. The two EGP-NPA replicates use $481990$ and $484821$ steps. Final-budget checkpoints are specified before testing, without selection on test results.

Native parameterizations and learning objectives are retained and disclosed. DGPPO receives 6-dimensional node and 5-dimensional edge features containing relative goals, relative geometry, boundary clearance, and exact spectral margin. EGP-NPA receives 17/20-dimensional inputs including its shield's internal state. DGPPO's task reward is normalized progress minus the time-step cost; it does not include EGP-NPA's dual-dependent stagnation term or terminal arrival reward. These differences are components of the complete systems, rather than a one-factor intervention isolating the policy-update formula.

DGPPO's policy, task-value network, and constraint-value network contain $37796$, $33441$, and $14211$ parameters, respectively; the EGP-NPA actor contains $108998$. Native network sizes are retained rather than parameter-matched, so observed differences also include representation and control-architecture effects.

Testing uses $96$ newly generated and frozen tasks: $64$ at $N=8$ and $12,12,8$ at $N=12,20,25$. Two methods and two training replicates give $384$ closed-loop runs with identical initial states, goals, safety distance, speed limit, time horizon, holding criterion, and independent rational auditing. DGPPO executes its policy directly without our shield. EGP-NPA's allocation and holding-segment certification are explicit system components; certification and computational-budget stops remain separate outcomes.

\begin{table*}[t]
\centering\footnotesize
\caption{Adapted External-Method Comparison on 96 New Tasks. Non-Safe-Arrivals Receive 60 Seconds in Mean Completion Time; Audit Failure Is Not Necessarily a Physical Collision.}\label{tab:external}
\begin{tabularx}{\textwidth}{@{}Xccccc@{}}\toprule
Method & Replicate & $N=8$ Safe & Cross-Scale Safe & Failed-Audit Runs & Mean Completion/s\\\midrule
EGP-NPA & 1 & 64/64 & 29/32 & 0/96 & 13.67 \\
EGP-NPA & 2 & 64/64 & 28/32 & 0/96 & 13.67 \\
DGPPO (task-adapted) & 1 & 0/64 & 0/32 & 96/96 & 60.00 \\
DGPPO (task-adapted) & 2 & 0/64 & 0/32 & 96/96 & 60.00 \\
\bottomrule\end{tabularx}
\end{table*}

For EGP-NPA replicate 1, nominal arrivals are $93/96$, safe arrivals are $93/96$, and $0$ runs have physical collisions.
For EGP-NPA replicate 2, nominal arrivals are $92/96$, safe arrivals are $92/96$, and $0$ runs have physical collisions.
For DGPPO replicate 1, nominal arrivals are $0/96$, safe arrivals are $0/96$, and $96$ runs have physical collisions.
For DGPPO replicate 2, nominal arrivals are $0/96$, safe arrivals are $0/96$, and $96$ runs have physical collisions.


Table~\ref{tab:external} evaluates these implementations at the stated budgets. Only two training replicates are available; finite-bank versus online sampling, training rewards, and prior tuning effort are not fully matched. The global spectral constraint is a task extension of DGPPO, so this result does not replace its published benchmark conclusions. We make neither algorithm-level significance nor general-superiority claims. Failure of an adapted policy to meet safety requirements describes the present configuration, not an impossibility result for discrete-barrier learning.
